\prefacesection{Abstract}

% TODO: revisit after the developer-interview component is complete (RQ6--RQ9).

Software for motion planning and simulation (MPS) supports the design, testing,
and deployment of robotic systems. Because experiments and applications depend
on this software behaving correctly, its engineering quality matters beyond any
single project. This report assesses the state of the practice for MPS software
by adapting the research-software assessment methodology of Smith et al.
Candidate packages were gathered from academic surveys, community-curated
lists, and auxiliary LLM-assisted discovery; after consolidation of 295 unique
candidate names and filtering by open-source availability, repository-based
measurability, and maintenance activity, 28 packages were selected, spanning
full simulators, physics and dynamics engines, motion planning libraries, and
middleware. Each package was measured in May 2026 against a standard template
covering installability, correctness and verifiability, surface reliability,
surface robustness, surface usability, maintainability, reusability, surface
understandability, and visibility and transparency, together with repository
metrics. A commonality analysis describes the selected packages as a software
family.

The measured packages show strong practical infrastructure: all 28 provide
installation instructions, installation automation, tutorials, user manuals,
and API documentation, and only two installations broke during measurement.
Weaknesses concentrate in formal artifacts: only one package documents expected
user characteristics, and half do not identify a coding standard. Measured
quality and community popularity diverge for several packages, so popularity
indicators such as GitHub stars are not reliable proxies for engineering
quality. Compared with prior studies of the state of the practice for medical imaging
and Lattice Boltzmann method software, MPS software is at least as well
equipped in repository artifacts and development tooling. A developer-interview
component has been designed to investigate pain points and practices; its
execution is left for future work.
